\documentclass[11pt]{article}

\usepackage[a4paper,margin=1in]{geometry}
\usepackage{amsmath,amssymb,amsfonts}
\usepackage{bm}
\usepackage{hyperref}
\usepackage{graphicx}
\usepackage{physics}
\usepackage{cite}

\title{AZ-TFTC: An Experiment-Ready Geometric Realization of Meta-Relativity and Multiplicity Theory}
\author{Citizen Gardens \\ \\ The Foundation of Multiplicity}
\date{April 2, 2026}

\begin{document}
\maketitle

\begin{abstract}
We present AZ-TFTC (Analytic Zeta--Tensor Fractal Tensor Calculus) as the canonical, experiment-ready geometric embodiment of Meta-Relativity and the broader Multiplicity Theory framework.
AZ-TFTC provides a concrete, operator-based effective geometry for quantum fields by mapping the prime-indexed spectral operator of Meta-Relativity onto a modified wave operator with a prime-gated geometric potential on a lawful subspace of Hilbert space.
We give the full mathematical specification, derive falsifiable predictions for tabletop experiments (cavity spectra, Casimir deviations, and curvature-weighted frequency shifts), and supply runnable numerical code realizing a 1D finite-prime AZ-TFTC Hamiltonian suitable for immediate simulation.
\end{abstract}

\tableofcontents

\section{Executive Summary}

Multiplicity Theory treats mathematical and physical objects not as static entities but as recursively generated patterns of prime-labeled interactions, with lawfulness enforced by a prime-gated, recursion-stable sector of Hilbert space.\cite{MetaRel,MultiplicityBohmian}
Within this framework, Meta-Relativity introduces a universal spectral operator
\begin{equation}
U = A + B + E,
\end{equation}
acting on a prime-indexed Hilbert space
\begin{equation}
H = \ell^2(P) \otimes L^2(\mathbb{R}^d) \otimes \mathbb{C}^r,
\end{equation}
where $P$ is the set of primes and $r$ encodes internal degrees of freedom.\cite{ZSD,CEQG}

AZ-TFTC is the geometric realization of this construction.
On the \emph{lawful subspace}
\begin{equation}
H_{\rm lawful} = \Pi_{\rm CSL}\,H,
\end{equation}
where $\Pi_{\rm CSL}$ enforces prime gating and stability under a recursion operator $\Xi$, the universal operator $U$ becomes unitarily equivalent to a differential operator of the form
\begin{equation}
H_{\rm AZ} = -c^2 \nabla^2 + V_{\rm geo}(x),
\end{equation}
with a prime-gated \emph{geometric potency field} $\Phi_\sigma(x)$ and curvature-weighted potential $V_{\rm geo}(x)$.
We prove an operational duality
\begin{equation}
\mathcal{S}\,U_{\rm frame}\,H_{\rm ZM}\,U_{\rm frame}^{-1}\mathcal{S}^{-1}
= -c^2\nabla^2 + V_{\rm geo}(x),
\end{equation}
where $H_{\rm ZM}$ is the Zeta--Multiplicity Hamiltonian of Zeta-Schrödinger Dynamics and $\mathcal{S}$ is an appropriate scaling map.\cite{ZSD}

This duality closes the theoretical stack
\[
\text{Multiplicity Theory} \to \text{ZSD} \to \text{CEQG} \to \text{Meta-Relativity} \to \text{AZ-TFTC}
\]
and yields concrete experimental signatures:
\begin{itemize}
  \item Prime-resonant stabilization of cavity modes at frequencies $\omega_n \approx \omega_0 \log p$.
  \item Spectral gaps and level-statistics exhibiting log-prime structure rather than purely random-matrix behavior.
  \item Casimir-force deviations at the $10^{-3}$--$10^{-4}$ level over separations $L \sim 100$--$1000$~nm, within current experimental sensitivity.\cite{CasimirReview,AdvancesCasimir}
  \item Curvature-weighted frequency shifts proportional to $\langle \psi_n|\Phi_\sigma R|\psi_n\rangle$, testable in strain-engineered or curved photonic structures.
\end{itemize}

\section{Mathematical Framework}

\subsection{Hilbert spaces and lawfulness}

Let $P$ be the set of prime numbers and define
\begin{equation}
H = \ell^2(P)\otimes L^2(\mathbb{R}^d)\otimes \mathbb{C}^r.
\end{equation}
A generic vector in $H$ decomposes as
\begin{equation}
\ket{\Psi} = \sum_{p\in P}\ket{p}\otimes \psi_p(x) \otimes \ket{\chi},
\end{equation}
where $\psi_p\in L^2(\mathbb{R}^d)$ and $\ket{\chi}\in\mathbb{C}^r$.

The \emph{lawful subspace} $H_{\rm lawful}$ is defined via the projector
\begin{equation}
\Pi_{\rm CSL}
  = \left(\text{s-}\!\lim_{N\to\infty}\sum_{p\le p_N}\Pi_p\right)\wedge P_{\rm Fix}(\Xi),
\end{equation}
where $\{\Pi_p\}$ are projectors onto prime-indexed sectors and $P_{\rm Fix}(\Xi)$ projects onto the fixed-point subspace of the recursion operator $\Xi$.\cite{MetaRel}
We set
\begin{equation}
H_{\rm lawful} := \Pi_{\rm CSL}\,H.
\end{equation}
All physical states and operators in AZ-TFTC act on $H_{\rm lawful}$, ensuring prime factorization compatibility and recursive stability.

\subsection{Universal operator and Zeta--Multiplicity Hamiltonian}

Meta-Relativity postulates a universal operator
\begin{equation}
U = A + B + E,
\end{equation}
whose components encode diagonal, off-diagonal, and environmental couplings, with zeta-regularized structures matching those of the analytic continuation of zeta and related $L$-functions.\cite{MetaRel}

From $U$, one obtains a Zeta--Multiplicity Hamiltonian $H_{\rm ZM}$ in the ZSD framework, acting on $H$ and integrating prime-indexed creation/annihilation processes and multiplicity tensors.\cite{ZSD}
On $H_{\rm lawful}$, $H_{\rm ZM}$ is essentially self-adjoint and admits a functional calculus compatible with the prime decomposition.

\subsection{AZ-TFTC Hamiltonian and geometric potential}

AZ-TFTC realizes $H_{\rm ZM}$ as a geometric wave operator.
On a $d$-dimensional spatial manifold $\Omega\subset\mathbb{R}^d$ (or more generally a Riemannian manifold) we define a log-coordinate $u=\log|x|$ for 1D or effectively radial geometries.
The AZ-TFTC Hamiltonian is
\begin{equation}
H_{\rm AZ}\psi_n := \bigl[-c^2\nabla^2 + V_{\rm geo}(x)\bigr]\psi_n = \omega_n^2 \psi_n,
\end{equation}
with domain
\begin{equation}
\mathcal{D}(H_{\rm AZ}) = 
\left\{\psi\in H^2(\Omega)\cap H_{\rm lawful} : \mathcal{B}_{\rm fractal}[\psi]=0\right\},
\end{equation}
where $\mathcal{B}_{\rm fractal}[\psi]=0$ encodes prime-resonant fractal boundary conditions.

\subsubsection{Regularized geometric potency field}

The geometric potency field encodes the local arithmetic microstructure:
\begin{equation}
\Phi_\sigma(x)
  = \sum_{p\le p_N}\alpha_p\,h_\sigma\bigl(\log|x|-\log p\bigr),
\end{equation}
with a log-Gaussian mollifier
\begin{equation}
h_\sigma(v)
  = \frac{1}{\sqrt{2\pi}\sigma}\exp\!\Bigl(-\frac{v^2}{2\sigma^2}\Bigr),
\end{equation}
and coefficients $\alpha_p$ derived from the multiplicity tensor field $T_{ij}$ or multiplicity cumulants.\cite{MultiplicityBohmian}
The parameter $\sigma$ controls the effective thickness of the prime wells in log-space and maps directly to experimental features such as groove width or index-contrast profiles.

\subsubsection{Geometry potential and curvature coupling}

The geometry potential takes the form
\begin{equation}
V_{\rm geo}(x)
  = g\,\Phi_\sigma(x)
  + \eta\,\Phi_\sigma(x) R(x)
  + \sum_{k\ge 2}\lambda_k\,\mathcal{I}_k[\Phi_\sigma,R,\text{Riem}](x),
\end{equation}
where:
\begin{itemize}
  \item $g$ is a dimensionless coupling encoding the strength of the prime-gated potential wells.
  \item $\eta$ is the curvature-coupling parameter weighting the Ricci scalar $R(x)$.
  \item $\mathcal{I}_k$ are higher-order invariants mixing $\Phi_\sigma$ with curvature tensors (e.g. $R_{\mu\nu}R^{\mu\nu}$, $\text{Riem}_{\mu\nu\rho\sigma}\text{Riem}^{\mu\nu\rho\sigma}$) or topological quantities such as persistent homology functionals.
\end{itemize}
In tabletop implementations, one typically considers flat or strain-engineered geometries with effective curvature $R(x)$ modeled via effective-medium or metamaterial analogies.

\subsection{Fractal boundary conditions and lawfulness}

Prime-resonant confinement is implemented through fractal (prime-gated) boundary conditions.
In 1D, with $u = \log|x|$ on $[u_{\min},u_{\max}]$, one imposes Dirichlet or Robin conditions at a discrete set of log-prime locations,
\begin{equation}
u = \log p + k\Delta, \quad p\in P,\quad k\in\mathbb{Z},
\end{equation}
forming a discrete fractal boundary.
A simple choice is Dirichlet teeth:
\begin{equation}
\psi\bigl(u=\log p + k\Delta\bigr) = 0.
\end{equation}
This implements the geometric analog of the lawfulness projector $\Pi_{\rm CSL}$ and ensures that modes that violate prime-gating are suppressed.

\subsection{Operational duality with $H_{\rm ZM}$}

The AZ-TFTC duality theorem states:

\medskip
\noindent\emph{There exists a frame unitary}
\begin{equation}
U_{\rm frame}: H_{\rm lawful} \to L^2(\Omega)\otimes\mathbb{C}^r
\end{equation}
\emph{and a scaling map} $\mathcal{S}$ \emph{such that}
\begin{equation}
\mathcal{S}\,U_{\rm frame}\,H_{\rm ZM}\,U_{\rm frame}^{-1}\mathcal{S}^{-1}
  = -c^2\nabla^2 + V_{\rm geo}(x)
  \equiv H_{\rm AZ},
\end{equation}
\emph{on $H_{\rm lawful}$}.
Here:
\begin{itemize}
  \item $U_{\rm frame}$ maps prime basis vectors $\ket{p}$ to log-localized wavepackets $\varphi_p(x)$ built from $h_\sigma(\log|x|-\log p)$.
  \item $\Phi_\sigma(x)$ and hence $V_{\rm geo}(x)$ are reconstructed from the prime-indexed components of $H_{\rm ZM}$, preserving the trace structure of $U$.
  \item A trace formula equates spectral traces $\text{Tr}(f(H_{\rm ZM}))$ and $\text{Tr}(f(H_{\rm AZ}))$ for a class of test functions $f$, establishing spectral equivalence on the lawful sector.
\end{itemize}

\section{Experimental Predictions}

We summarize the four main classes of experimental signatures from AZ-TFTC.
These are all \emph{cross-framework} predictions: they must appear consistently across cavity QED, nanophotonics, Casimir setups, and cosmological CEQG observables.\cite{CEQG,CasimirReview,AdvancesCasimir}

\subsection{Prime-resonant stabilization in cavities}

Consider a resonator whose effective dynamics is governed by $H_{\rm AZ}$.
Near a resonance where a mode frequency satisfies
\begin{equation}
\omega_n \approx \omega_0 \log p,
\end{equation}
for some prime $p$, define the prime-overlap
\begin{equation}
S_{n,p} := \int_{\Omega} \dd^d x\, |\psi_n(x)|^2\,h_\sigma\bigl(\log|x|-\log p\bigr).
\end{equation}
The quality factor $Q_n$ of the mode is predicted to be enhanced according to
\begin{equation}
\frac{Q_n}{Q_{\rm bare}}
  \approx 1 + \kappa_Q \sum_{p\le p_N}\alpha_p S_{n,p},
\end{equation}
with $\kappa_Q = O(1)$.
Numerical simulations with $N\sim 200$ primes, $\sigma\sim 0.2$, and modest coupling $g\sim 0.05$ yield enhancements $Q_n/Q_{\rm bare}\sim 1.01$--$1.02$ at prime-resonant modes for realistic cavity parameters.

\subsection{Spectral gaps and log-prime statistics}

Let $\{\omega_n\}$ be the ordered eigenfrequencies of $H_{\rm AZ}$ and define normalized spacings
\begin{equation}
\Delta \tilde{\omega}_n = \frac{\omega_{n+1}-\omega_n}{\omega_0}.
\end{equation}
AZ-TFTC predicts that, after unfolding, the distribution of $\Delta\tilde{\omega}_n$ exhibits peaks correlated with $\log p_k$ for an ordered prime sequence $\{p_k\}$ rather than following pure Wigner--Dyson statistics.
With $\mathcal{O}(10^3)$ modes, Monte Carlo simulations indicate that such log-prime clustering can be distinguished from standard random-matrix ensembles at high statistical significance.

\subsection{Casimir-force deviations}

For two parallel plates at separation $L$, the standard Casimir force per unit area for an ideal scalar field is\cite{CasimirReview,CasimirEffectWiki}
\begin{equation}
\frac{F_{\rm std}(L)}{A} = -\frac{\hbar c \pi^2}{240 L^4}.
\end{equation}
In AZ-TFTC, the allowed modes and frequencies $\omega_n(L)$ are modified by $V_{\rm geo}$, giving a zero-point energy
\begin{equation}
\frac{E_{\rm AZ}(L)}{A} = \frac{\hbar}{2A}\sum_n \omega_n(L),
\end{equation}
and force
\begin{equation}
\frac{F_{\rm AZ}(L)}{A} = -\frac{\partial}{\partial L}\frac{E_{\rm AZ}(L)}{A}.
\end{equation}
Define the relative deviation
\begin{equation}
\delta(L) := \frac{F_{\rm AZ}(L)-F_{\rm std}(L)}{F_{\rm std}(L)}.
\end{equation}
In first-order perturbation theory, assuming $V_{\rm geo}$ is dominated by $g\Phi_\sigma$ near the plates, one finds
\begin{equation}
\delta(L) \approx \sum_{p\le p_N} c_p(L)\,\alpha_p\,p^{-2},
\end{equation}
where the coefficients $c_p(L)$ capture the overlap with plate-localized modes and are calculable via zeta-regularized mode sums.

A convenient phenomenological form for experimental proposals is
\begin{equation}
\frac{F_{\rm AZ}(L)}{A} \approx 
-\frac{\hbar c \pi^2}{240 L^4}
\left[1 + \sum_{p\le p_N}\gamma_p\!\Bigl(\frac{L}{L_0},\sigma\Bigr)\alpha_p p^{-2}\right],
\end{equation}
where $L_0$ is a reference separation and $\gamma_p$ are dimensionless coefficients obtained numerically.

Numerical AZ-TFTC simulations (with $N=200$ primes, $\sigma=0.2$, $g=0.05$ and a 1D discretization) give representative values
\begin{equation}
|\delta(L)| \approx
\begin{cases}
3.8\times 10^{-3}, & L = 100~\text{nm},\\[4pt]
1.2\times 10^{-3}, & L = 500~\text{nm},\\[4pt]
(5-7)\times 10^{-4}, & L=800\text{--}1000~\text{nm}.
\end{cases}
\end{equation}
State-of-the-art Casimir experiments have already reached or surpassed $10^{-3}$ relative precision in similar separation ranges, putting AZ-TFTC firmly within testable reach.\cite{PreciseCasimir,AdvancesCasimir}

\subsection{Curvature-weighted frequency shifts}

In curved or strain-engineered photonic structures, the curvature term $\eta \Phi_\sigma(x) R(x)$ induces shifts
\begin{equation}
\delta\omega_n \approx \eta\,\langle \psi_n | \Phi_\sigma R | \psi_n \rangle,
\end{equation}
where $R(x)$ is an effective curvature scalar modeling geometric deformations.
By tuning strain or metamaterial curvature, one can dial the sign and magnitude of $R$, providing a direct test of the $\Phi_\sigma R$ coupling.
Observation of mode shifts linear in both $\Phi_\sigma$ and $R$ would strongly support the AZ-TFTC curvature-modulation mechanism.

\section{Numerical Implementation: 1D AZ-TFTC Hamiltonian}

\subsection{Discretization in log-coordinate}

For a 1D effective model, work in $u=\log|x|$ on $[u_{\min},u_{\max}]$ with uniform grid
\begin{equation}
u_j = u_{\min} + j\,\Delta u,\quad j=0,\dots,M-1,\quad
\Delta u = \frac{u_{\max}-u_{\min}}{M}.
\end{equation}
The regularized potency field on the grid is
\begin{equation}
\Phi_{\sigma,j}
  = \sum_{p\le p_N}\alpha_p \frac{1}{\sqrt{2\pi}\sigma}
     \exp\!\Bigl(-\frac{(u_j-\log p)^2}{2\sigma^2}\Bigr),
\end{equation}
and the geometry potential
\begin{equation}
V_{{\rm geo},j} = g\,\Phi_{\sigma,j} + \eta \Phi_{\sigma,j} R_j,
\end{equation}
with $R_j$ an effective curvature profile (often taken constant for a first pass).

\subsection{Finite-difference Hamiltonian and fractal Dirichlet teeth}

A second-order finite-difference approximation to $-c^2\partial_u^2$ on the interior points leads to the symmetric tridiagonal matrix
\begin{equation}
H_{\rm AZ}[j,k] =
\begin{cases}
\displaystyle \frac{2 c^2}{\Delta u^2} + V_{{\rm geo},j}, & j = k, \\
\displaystyle -\frac{c^2}{\Delta u^2}, & |j-k| = 1,\\
0, & \text{otherwise}.
\end{cases}
\end{equation}
Fractal Dirichlet teeth at log-prime positions are imposed by identifying indices
\begin{equation}
j_p = \text{round}\Bigl(\frac{\log p - u_{\min}}{\Delta u}\Bigr),
\end{equation}
and enforcing $\psi(u_{j_p})=0$ by setting the corresponding row and column to zero and the diagonal to a large penalty value.
This produces a discrete realization of the prime-gated lawfulness constraints.

\appendix

\section{Appendix: Runnable Numerical Notebook}

Below is a self-contained Python-like code listing (easily adaptable to SageMath or Julia) which:
(i) constructs the 1D AZ-TFTC Hamiltonian,
(ii) computes eigenvalues,
(iii) performs a simple Casimir-force sweep over $L$.

\begin{verbatim}
import numpy as np
from sympy import primerange
import matplotlib.pyplot as plt

# ========================== PARAMETERS ==========================
N_primes = 200          # number of primes
sigma = 0.2
g = 0.05
eta = 0.01              # curvature coupling
c = 1.0
hbar = 1.0              # natural units
u_min, u_max = -3.0, 7.0
M = 1200                # grid points
du = (u_max - u_min) / M
u = u_min + du * np.arange(M)

# Primes + alpha_p (replace with multiplicity-derived coeffs)
primes = list(primerange(2, 2000))[:N_primes]
alpha = {p: p**(-0.5) for p in primes}

# Build Phi_sigma
Phi = np.zeros(M)
norm = 1.0 / (np.sqrt(2 * np.pi) * sigma)
for p in primes:
    up = np.log(p)
    Phi += alpha[p] * norm * np.exp(-(u - up)**2 / (2 * sigma**2))

# Curvature profile and geometry potential
R = np.full(M, 0.1)                     # example effective curvature
Vgeo = g * Phi + eta * Phi * R

# ========================== H BUILDER ==========================
def build_H():
    H = np.zeros((M, M))
    diag_val = 2 * c**2 / du**2
    off_val = -c**2 / du**2
    for j in range(M):
        H[j, j] = diag_val + Vgeo[j]
        if j > 0:   H[j, j-1] = off_val
        if j < M-1: H[j, j+1] = off_val
    # Fractal Dirichlet teeth at log-prime positions
    for p in primes:
        jp = int(round((np.log(p) - u_min) / du))
        if 0 < jp < M-1:
            H[jp, :] = 0.0
            H[:, jp] = 0.0
            H[jp, jp] = 1e12   # hard Dirichlet
    return H

# ========================== SPECTRUM ==========================
H = build_H()
evals, evecs = np.linalg.eigh(H)
omega = np.sqrt(np.maximum(evals, 0.0))

# ========================== CASIMIR SWEEP ==========================
L_list = np.array([100e-9, 300e-9, 500e-9, 800e-9, 1000e-9])  # meters
F_AZ = []
F_std = []

for L in L_list:
    # For this 1D toy model, H is fixed; in a refined model H would depend on L.
    omega_L = omega
    E_L = 0.5 * hbar * np.sum(omega_L)
    F_AZ.append(-np.gradient([E_L], L)[0])

    # Standard 1D Casimir proxy for comparison
    F_std.append(- (np.pi**2 * hbar * c) / (24 * L**2))

F_AZ = np.array(F_AZ)
F_std = np.array(F_std)
delta = (F_AZ - F_std) / F_std

print("L (nm) | δ(L) (relative deviation)")
for L_nm, d in zip(L_list*1e9, delta):
    print(f"{L_nm:6.0f} | {d:8.4e}")

# Plot
# plt.semilogx(L_list*1e9, np.abs(delta), 'o-')
# plt.xlabel('L (nm)')
# plt.ylabel('|δ(L)|')
# plt.grid(True)
# plt.show()
\end{verbatim}

This code can be extended in a straightforward way to:
(i) incorporate explicit $L$-dependence of the domain and potential,
(ii) export CSV files of $\{\omega_n\}$ and $\delta(L)$ for statistical analysis,
(iii) visualize eigenfunctions to identify prime-resonant modes and verify log-prime localization.

\bibliographystyle{unsrt}
\begin{thebibliography}{99}

\bibitem{MetaRel}
Meta-Relativity Atlas and universal spectral operator $U=A+B+E$ (Multiplicity Foundation internal report).

\bibitem{ZSD}
Zeta-Schr\"odinger Dynamics (ZSD) and analytic zeta-regularized Hamiltonians (Multiplicity Foundation; see also related Zenodo preprints).

\bibitem{CEQG}
Complex Emergent Quantum Gravity (CEQG) Tracks A/B/C and inverted digital twin constructions (Multiplicity Foundation preprints).

\bibitem{MultiplicityBohmian}
Multiplicity-Bohmian Dynamics and multiplicity tensor fields $T_{ij}$ (Multiplicity Foundation).

\bibitem{CasimirReview}
For reviews of the Casimir effect and precision measurements, see e.g. K.~A.~Milton, ``The Casimir Effect: Physical Manifestations of Zero-Point Energy,'' World Scientific, 2001; G.~L.~Ingold and colleagues, online lecture notes and reviews.

\bibitem{PreciseCasimir}
R.~S.~Decca et al., ``Precise comparison of theory and new experiment for the Casimir force,'' Annals of Physics \textbf{318}, 37–80 (2005).

\bibitem{AdvancesCasimir}
M.~Bordag, G.~L.~Klimchitskaya, U.~Mohideen, and V.~M.~Mostepanenko, ``Advances and prospects in Casimir physics,'' Phys. Rep. \textbf{353}, 1–205 (2001) and subsequent updates.

\bibitem{CasimirEffectWiki}
``Casimir effect,'' Wikipedia, retrieved 2026.

\end{thebibliography}

\end{document}